\documentclass[letterpaper]{article}
\PassOptionsToPackage{unicode=true}{hyperref} % options for packages loaded elsewhere
\PassOptionsToPackage{hyphens}{url}
\usepackage{graphicx} % Required for inserting images


\def\tightlist{}


\graphicspath{ {../images/} }
\usepackage[margin=1in]{geometry}
\usepackage{tikz}
\usepackage[dvipsnames]{xcolor}
\usetikzlibrary{calc}




\usepackage[utf8]{inputenc}
\usepackage[T1]{fontenc}
\usepackage{lmodern}
\usepackage[english]{babel}

\usepackage{longtable}
\usepackage{booktabs}

% Useful packages
\usepackage{hyperref}
\usepackage{amsmath, amssymb}
\usepackage{fancyhdr}
\usepackage{setspace}
\usepackage{tocloft}
\usepackage{titlesec}


\usepackage{unicode-math}



% Code highlighting (Pandoc uses listings or minted, but listings is safer without shell escape)
\usepackage{listings}
\lstset{
  basicstyle=\ttfamily\small,
  breaklines=true,
  frame=single,
  backgroundcolor=\color[gray]{0.95}
}

\usepackage{setspace}


%Modifiable Commands; defaults set to none
\newcommand{\titleDOC}{}
\newcommand{\authorDOC}{Daniel Sabalakov}
\newcommand{\dateDOC}{\today}
\newcommand{\classDOC}{}
\newcommand{\emailDOC}{danielsabalakov@gmail.com}
% DEFAULT QUOTE IS BY ISSAC NEWTON. IF YOU DON'T WANT IT, DO: quote: @none
\newcommand{\quoteDOC}{``\textit{Truth is ever to be found in the simplicity, and not in the multiplicity and confusion of things}'' - Issac Newton}

% Header/Footer
\pagestyle{fancy}
\fancyhf{}
\rhead{\thepage}
\lhead{\leftmark}





% %Title Arguments
% \title{\TITLEOFDOCUMENT}
% \author{\AUTHOROFDOCUMENT}
% \date{\today}


%
% ANYTHING BELOW HERE CAN GET PROCEDURALLY GENERATED
%
% BEGINNING OF DOCUMENT
%
%
%
\begin{document}

%titlepage








\renewcommand{\authorDOC}{Daniel Sabalakov}

\renewcommand{\classDOC}{DE Intro to Biology I}

\renewcommand{\titleDOC}{Bubble Lab Exploration}





\begin{titlepage}
  \titleDOC
  \classDOC
  \dateDOC
  \authorDOC
\end{titlepage}

\newpage

\section{Background}\label{background}

Bubbles make a great stand in for cell membranes. They're fluid, flexible, and can self-repair. Bubbles and cell membranes are alike because their parts are so similar. If you could zoom down on a cell membrane, you'd see that much of the membrane is a double layer of little molecules called phospholipids.

Phospholipids have a love-hate relationship with water. One end, the ``head,'' is attracted to water (\textbf{hydrophilic}), and the other end, the ``tail,'' is repelled by water (\textbf{hydrophobic}). This dual nature is described as \textbf{amphipathic}. Place phospholipids in water and they quickly form a double layer with the heads facing out on both sides.

\textbf{A soap molecule has the same split personality. The ``head'' of a soap molecule is charged (ionic) and is attracted to water molecules, which have regions of positive and negative charge (polar). The hydrocarbon tail of the soap molecule is not charged and is repelled by water's polarity.} This explains why we use soap to clean. The hydrocarbon tail of soap mixes with and dissolves in other hydrocarbons, like oils and fats, while the head region grabs a hold of passing water molecules and follows them down the drain. The surface of a bubble has three layers. The middle layer is a thin film of water. On both sides of this film is a layer of soap molecules with hydrophilic heads oriented toward the water film and hydrophobic tails pointing away.

\subsection*{Analysis \& Questions}\label{analysis-questions}
\addcontentsline{toc}{subsection}{Analysis \& Questions}

In this lab, soap bubbles were used to model several properties that are characteristic of cellular membranes. Below is a chart listing each of the ``Cell Concepts'' investigated. For each concept:

\begin{enumerate}
\def\labelenumi{\arabic{enumi}.}
\item
  \textbf{Insert a photo or sketch of your group modeling this property}
\item
  \textbf{Explain. Follow the guiding questions to describe the cell concept, as you understand it, and how the soap bubble was used to model the cell concept.}
\item
  \textbf{Answer the questions that relate to that concept. You may use your reading guide or another source to answer this. Citations can be added as links at the end of the document.}
\end{enumerate}

\subsubsection*{Cell Concept 1 -- Fluid Mosaic Model: Membranes are Fluid and Flexible}\label{cell-concept-1-fluid-mosaic-model-membranes-are-fluid-and-flexible}
\addcontentsline{toc}{subsubsection}{Cell Concept 1 -- Fluid Mosaic Model: Membranes are Fluid and Flexible}

Describe how the bubble fits this description - Fluid? Mosaic? What is the ``mosaic'' in the cell?\\
What role do the following play in membrane fluidity? 1 - Temperature 2 - Fatty acid composition 3 - Cholesterol

\subsubsection*{Cell Concept 2 -- Membranes Can Self-Repair}\label{cell-concept-2-membranes-can-self-repair}
\addcontentsline{toc}{subsubsection}{Cell Concept 2 -- Membranes Can Self-Repair}

The bubble reformed after you removed your finger? What rejoins in a cell to repair small membrane tears?\\
Explain \textbf{why} the membrane wants to seek this repaired arrangement? What chemical properties are at play?

\subsubsection*{Cell Concept 3 -- Membrane Bound Organelles}\label{cell-concept-3-membrane-bound-organelles}
\addcontentsline{toc}{subsubsection}{Cell Concept 3 -- Membrane Bound Organelles}

What are some organelles in a cell that have membranes to define the compartment?\\
Explain why compartmentalization (separation by membranes within a cell) is important?

\subsubsection*{Cell Concept 4 -- Membrane Proteins}\label{cell-concept-4-membrane-proteins}
\addcontentsline{toc}{subsubsection}{Cell Concept 4 -- Membrane Proteins}

Membrane proteins are important for allowing for cell transport across the membrane. How did your thread model this? Do membrane proteins stay in one place / static?\\
Explain the structure and placement of a protein within the membrane? (think about the amphipathic nature of the phospholipids)

\subsubsection*{Cell Concept 5 -- Cell Connections}\label{cell-concept-5-cell-connections}
\addcontentsline{toc}{subsubsection}{Cell Concept 5 -- Cell Connections}

Why do you think cells might benefit from having a direct connection?\\
Explain and give examples of connections between cells like that modeled in the activity.

\subsubsection*{Cell Concept 6 -- Cell Division - Binary Fission / Cytokinesis}\label{cell-concept-6-cell-division---binary-fission-cytokinesis}
\addcontentsline{toc}{subsubsection}{Cell Concept 6 -- Cell Division - Binary Fission / Cytokinesis}

Why do cells need to divide? How does the division of cytoplasm in a plant differ from an animal cell to produce 2 cells from 1? Explain what happens to the cell membrane.
* Cells need to divide to replace old cells. After a cell begins dividing, it will eventually have two copies of the nucleus and other organelles. Then, the cellular membrane will regenerate as cytokinesis happens. In a cell wall, the plant cell just builds a cell wall between the old and new cell, but since an animal cell does not have that rigid structure, it needs to split completely via cytokinesis.

\subsection*{Sources?}\label{sources}
\addcontentsline{toc}{subsection}{Sources?}

Did you use any \textbf{sources}? If so, please include them here: (hyperlinks are fine)

\newpage



\end{document}
